\Subhead{4.5.1}{Молекулярный генератор. Квантовые стандарты частоты (времени)}

Молекулярный генератор использует микроволновый переход молекулы с предварительным отбором возбуждённых состояний. В аммиачном мазере инверсионное расщепление связано с двумя конфигурациями молекулы NH$_3$. Неоднородное электрическое поле сортирует состояния по эффекту Штарка; отобранный пучок проходит резонатор. Вынужденное излучение на переходе усиливает поле, если приток возбуждённых молекул превышает потери и декогеренцию.

Качественно уравнения поля и коллективной поляризации имеют вид
\begin{equation}
 \dot a=-\kappa_a a+gP,\qquad
 \dot P=-(\gamma_\perp+\ii\Delta)P+gNa.
\end{equation}
При $\Delta=0$ порог задаёт $g^2N>\kappa_a\gamma_\perp$. Частота формируется совместно переходом и резонатором; затягивание к резонатору, конечное время пролёта и столкновения сдвигают её от идеального перехода. Здесь $\kappa_a$ — амплитудный декремент, а не скорость потери энергии.

В пассивном стандарте атомы не обязаны самостоятельно генерировать. Стабилизируемый генератор облучает атомный ансамбль; по отклику формируют сигнал ошибки и подстраивают генератор к центру резонанса. Разделённые зоны Рамзе дают узкие интерференционные полосы с масштабом ширины $1/T$, где $T$ — свободное время между взаимодействиями. Узкая линия полезна лишь при достаточном отношении сигнал/шум.

Секунда СИ связана с невозмущённым сверхтонким переходом основного состояния $^{133}$Cs: $\Delta\nu_{\rm Cs}=9\,192\,631\,770$ Гц точно. Реальное устройство должно оценивать сдвиги Зеемана, Штарка, теплового излучения, столкновений, доплеровский и гравитационный. Стабильность обычно характеризуют относительными флуктуациями частоты на заданном времени, например девиацией Аллана; точность требует учёта систематических сдвигов. Активный водородный мазер и пассивный цезиевый стандарт играют разные роли, хотя оба основаны на квантовых переходах.
